\chapter{Evidence and Tracing}
\label{ch:evidence-tracing}

This chapter describes the evidence capture model used by \texttt{gert.home}: what is
captured, when, how it is stored, and how to verify and read the evidence bundle from a
completed run. The underlying trace format is defined by gert core (\emph{gert Design
Document}, §12); this chapter focuses on the \texttt{gert.home}-layer semantics.

% ─────────────────────────────────────────────────────────────────────────────
\section{The Evidence Model}
\label{sec:evidence-model}
% ─────────────────────────────────────────────────────────────────────────────

\texttt{gert.home} captures evidence at the step level. Each piece of evidence is associated
with the step that required it, the executor who submitted it, and the timestamp of submission.

\paragraph{What gets captured.}
\begin{itemize}
  \item \textbf{Photos} --- Binary image file stored under the run directory.
        Captured: filename, timestamp, zone and asset context, SHA-256 hash.
  \item \textbf{Notes} --- Free-text string submitted by the executor.
        Captured: note text, executor identity (name or delegate name), timestamp.
  \item \textbf{Checklists} --- Completion status for each item in the checklist.
        Captured: item label, checked/unchecked, timestamp per item, overall completion.
\end{itemize}

\paragraph{What is NOT evidence.}
The following events are recorded in the trace log but are \emph{not} included in the
evidence bundle:
\begin{itemize}
  \item Reminder notifications
  \item Overdue escalation alerts
  \item Scheduling events (\texttt{routine\_due}, \texttt{routine\_skipped})
  \item Delegation activation events
\end{itemize}

% ─────────────────────────────────────────────────────────────────────────────
\section{Evidence Record Structure}
\label{sec:evidence-record}
% ─────────────────────────────────────────────────────────────────────────────

Each evidence submission is recorded as a structured object in the trace log.

\begin{minted}{yaml}
type: evidence_submitted
seq: 3
ts: "2025-04-28T09:15:00Z"
step_id: pool_clean
evidence_type: photo
label: Pool water and test strip
submitted_by: Carlos Méndez
role: delegate
content:
  filename: pool_clean_2025-04-28.jpg
  sha256: a3f2b19e4c7d8e6f1a2b3c4d5e6f7a8b9c0d1e2f3a4b5c6d7e8f9a0b1c2d3e4f5
\end{minted}

\begin{description}
  \item[\texttt{type}] Always \texttt{evidence\_submitted} for evidence records.
  \item[\texttt{seq}] Sequence number within the run's trace log.
  \item[\texttt{ts}] ISO 8601 timestamp of submission.
  \item[\texttt{step\_id}] ID of the step that required this evidence.
  \item[\texttt{evidence\_type}] One of \texttt{photo}, \texttt{note}, \texttt{checklist}.
  \item[\texttt{label}] Human-readable label from the routine or incident step definition.
  \item[\texttt{submitted\_by}] Name of the person who submitted the evidence.
  \item[\texttt{role}] Either \texttt{homeowner} or \texttt{delegate}.
  \item[\texttt{content}] Type-specific payload. For photos: filename and SHA-256 hash.
        For notes: free text. For checklists: list of item completion records.
\end{description}

% ─────────────────────────────────────────────────────────────────────────────
\section{Run Directory Structure}
\label{sec:evidence-run-dir}
% ─────────────────────────────────────────────────────────────────────────────

Each routine run is stored in a directory under the configured run root
(\texttt{GERT\_HOME\_RUN\_DIR}, default: \texttt{./runs}).

\begin{minted}{yaml}
runs/
  pool_clean/
    2025-04-28/
      trace.jsonl          # Newline-delimited JSON trace events
      evidence/
        pool_clean_2025-04-28.jpg     # Photo evidence files
        pool_clean_2025-04-28.sha256  # SHA-256 sidecar file
    2025-05-05/
      trace.jsonl
      evidence/
        pool_clean_2025-05-05.jpg
        pool_clean_2025-05-05.sha256
    latest -> 2025-05-05   # Symlink to most recent run
  leak/
    2025-04-29/
      trace.jsonl
      evidence/
        assess_2025-04-29.jpg
        shutoff_2025-04-29.jpg
        fix_2025-04-29.jpg
\end{minted}

Each run date directory is named \texttt{YYYY-MM-DD} using the timezone from
\texttt{GERT\_HOME\_TZ}. The \texttt{latest} symlink always points to the most recent run
for quick access.

% ─────────────────────────────────────────────────────────────────────────────
\section{Reading the Timeline Projection}
\label{sec:evidence-timeline}
% ─────────────────────────────────────────────────────────────────────────────

The \texttt{trace.jsonl} file contains one JSON object per line, ordered by sequence number.
Key fields in each event:

\begin{description}
  \item[\texttt{seq}] Monotonically increasing sequence number within the run. Useful for
        ordering events when timestamps have the same second.
  \item[\texttt{ts}] ISO 8601 timestamp (UTC).
  \item[\texttt{step\_id}] ID of the step that emitted this event.
  \item[\texttt{type}] Event type. Home-layer types: \texttt{routine\_due},
        \texttt{routine\_completed}, \texttt{routine\_skipped}, \texttt{evidence\_submitted},
        \texttt{incident\_opened}, \texttt{incident\_step\_completed},
        \texttt{delegation\_active}.
\end{description}

A complete timeline for a pool-cleaning run looks like:

\begin{minted}{yaml}
{"seq":1,"ts":"2025-04-28T08:00:00Z","type":"delegation_active",
 "delegate_name":"Carlos Méndez","routine_id":"pool_clean"}
{"seq":2,"ts":"2025-04-28T08:00:01Z","type":"routine_due",
 "routine_id":"pool_clean","zone":"pool","executor":"Carlos Méndez"}
{"seq":3,"ts":"2025-04-28T09:15:00Z","type":"evidence_submitted",
 "step_id":"pool_clean","evidence_type":"photo",
 "submitted_by":"Carlos Méndez","sha256":"a3f2...c1d9"}
{"seq":4,"ts":"2025-04-28T09:15:01Z","type":"routine_completed",
 "routine_id":"pool_clean","duration_minutes":75}
\end{minted}

% ─────────────────────────────────────────────────────────────────────────────
\section{Evidence Bundles for Compliance}
\label{sec:evidence-bundles}
% ─────────────────────────────────────────────────────────────────────────────

For warranty claims, insurance documentation, or property sale due diligence, evidence can
be exported as a signed bundle. The bundle aggregates all evidence records for a routine or
date range.

\begin{minted}{yaml}
bundle:
  property: Casa Santiago
  generated_at: "2025-05-01T12:00:00Z"
  routine: pool_clean
  date_range:
    from: "2025-01-01"
    to: "2025-05-01"
  runs:
    - date: "2025-04-28"
      executor: Carlos Méndez
      role: delegate
      completed_at: "2025-04-28T09:15:01Z"
      evidence:
        - type: photo
          filename: pool_clean_2025-04-28.jpg
          sha256: a3f2b19e4c7d8e6f1a2b3c4d5e6f7a8b9c0d1e2f3a4b5c6d7e8f9a0b1c2d3e4f5
    - date: "2025-04-21"
      executor: homeowner
      role: homeowner
      completed_at: "2025-04-21T10:30:00Z"
      evidence:
        - type: photo
          filename: pool_clean_2025-04-21.jpg
          sha256: b4c3d2e1f0a9b8c7d6e5f4a3b2c1d0e9f8a7b6c5d4e3f2a1b0c9d8e7f6a5b4c3
\end{minted}

% ─────────────────────────────────────────────────────────────────────────────
\section{SHA-256 Verification}
\label{sec:evidence-sha256}
% ─────────────────────────────────────────────────────────────────────────────

Photo evidence files are SHA-256 hashed at the moment of submission. The hash is stored both
in the \texttt{trace.jsonl} event and in a \texttt{.sha256} sidecar file alongside the image.

To verify a photo has not been tampered with after submission:

\begin{minted}{yaml}
# The .sha256 sidecar contains the stored hash:
# a3f2b19e4c7d8e6f1a2b3c4d5...  pool_clean_2025-04-28.jpg

# Recompute and compare:
sha256sum pool_clean_2025-04-28.jpg
# If the output matches the sidecar, the file is unmodified.
\end{minted}

\begin{tcolorbox}[colback=yellow!10,colframe=orange!70,title=Note]
  SHA-256 verification proves the file has not been modified \emph{after} submission. It does
  not prove the photo was taken at the claimed location or time. For high-stakes compliance
  use cases, pair photo evidence with metadata (\texttt{EXIF} timestamp and GPS coordinates)
  captured at submission time.
\end{tcolorbox}

% ─────────────────────────────────────────────────────────────────────────────
\section{Reading Evidence from a Completed Run}
\label{sec:evidence-reading}
% ─────────────────────────────────────────────────────────────────────────────

To inspect the evidence from the most recent run of any routine, read the
\texttt{trace.jsonl} file under the \texttt{latest} symlink:

\begin{minted}{yaml}
# List all evidence events from the most recent pool_clean run:
# runs/pool_clean/latest/trace.jsonl

# Filter for evidence_submitted events:
{"seq":3,"ts":"2025-04-28T09:15:00Z","type":"evidence_submitted",
 "step_id":"pool_clean","evidence_type":"photo",
 "label":"Pool water and test strip",
 "submitted_by":"Carlos Méndez","role":"delegate",
 "content":{"filename":"pool_clean_2025-04-28.jpg",
             "sha256":"a3f2...c1d9"}}
\end{minted}

The evidence photos are stored at:
\texttt{runs/\{routine-id\}/\{date\}/evidence/\{filename\}}.

Each evidence file is self-contained: the filename, SHA-256 hash, and submitter identity are
all recorded in the trace, so the evidence bundle can be verified and read independently of
the gert runtime.
